%% =============================================================================
%% Section 1 — Introduction
%% "How Do LLM-Based Agents Represent Code Repositories?"
%% ACM Transactions on Software Engineering and Methodology (TOSEM)
%% =============================================================================

\section{Introduction}
\label{sec:introduction}

Modern software development is intrinsically a \emph{repository-level} activity.
A single developer task — fixing a bug, completing a function, resolving a GitHub
issue — requires reasoning across dozens of interdependent files, import graphs,
call chains, and test suites.
Yet the majority of research on applying Large Language Models (LLMs) to Software
Engineering (SE) treats code as an isolated snippet: feed a function to the model,
obtain a prediction.
This design choice is not a minor simplification.
It is a fundamental representational mismatch between the problems practitioners
face and the problems researchers solve.

The central question this paper addresses is therefore deceptively simple but
remarkably unexplored:
\begin{quote}
\emph{How do LLM-based systems actually represent a code repository when they need
to reason about it as a whole?}
\end{quote}

We answer this question through a rigorous Systematic Literature Review (SLR)
spanning 159 primary studies (2020--2026), producing a novel seven-paradigm
taxonomy of the \emph{Repository Representation Problem} (RRP) — the challenge
of selecting, encoding, and injecting the right subset of a large codebase into a
finite LLM context window — and mapping each paradigm against seven canonical SE
objectives: code completion, code generation, program repair, fault localization,
code search, code translation, and test generation.

\subsection{The Repository Representation Problem}
\label{subsec:rrp}

We formally define the RRP as follows.
Let $\mathcal{R}$ denote a repository consisting of $N$ source files
$\{f_1, \ldots, f_N\}$.
Given a natural-language or code query $q$, a \emph{repository representation
function} $\mathcal{F}$ maps $(\mathcal{R}, q)$ to a context $C$:
\[
  C = \mathcal{F}(\mathcal{R},\, q,\, \tau),
  \quad |C| \leq \tau,
\]
where $\tau$ is the LLM context-window budget (in tokens).
The function $\mathcal{F}$ must simultaneously optimise \emph{relevance}
(high signal-to-noise ratio), \emph{completeness} (capture all cross-file
dependencies), \emph{freshness} (reflect the current repository state), and
\emph{efficiency} (fit within $\tau$ at acceptable computational cost).
These four objectives are in tension; no existing approach optimises all four
simultaneously, which is precisely why seven distinct representation paradigms
have emerged in the literature.

\subsection{Motivation and Research Gap}
\label{subsec:motivation}

The last three years have seen an explosion of work at the intersection of LLMs
and software repositories.
Tools such as SWE-agent~\cite{yang2024sweagent}, AutoCodeRover
\cite{zhang2024autocoderover}, RepoCoder~\cite{zhang2023repocoder}, and
CodePlan~\cite{bairi2024codeplan} have demonstrated impressive capabilities on
repository-level benchmarks~\cite{jimenez2024swebench,liu2024repobench,
ding2023crosscodeeval}.
However, a unifying conceptual framework that explains \emph{why} these systems
make different design choices — and what the trade-offs are — has been
conspicuously absent.
Existing secondary studies, while valuable, do not fill this gap:

\begin{itemize}
  \item \textbf{Hou et al.}~\cite{hou2024llmse_slr} conduct a landmark SLR
        across \emph{all} LLM4SE tasks (395 papers, 2017--2024).
        Their breadth is a strength and a limitation: repository representation
        is one among dozens of topics, receiving no dedicated taxonomic
        treatment.
        The survey categorises \emph{what} LLMs do (code generation, repair,
        testing, \ldots) but not \emph{how} they ingest a repository to do so.

  \item \textbf{Husein et al.}~\cite{husein2025llmcompletion_slr} focus
        specifically on LLM-based code completion, a direct antecedent of our
        work.
        Their scope, however, is \emph{function-level} and
        \emph{file-level} completion; cross-file and repository-level
        representation strategies are outside their inclusion criteria.
        Their study predates the agentic paradigm and the RAG revolution in
        software engineering.

  \item \textbf{Jiang et al.}~\cite{jiang2026surveycodegen} survey LLMs for
        code generation with admirable depth (published in ACM TOSEM).
        Their taxonomy organises papers by \emph{generation target} (NL-to-code,
        code-to-code, test generation).
        The fundamental question of how a repository is \emph{represented} as
        context — the input side of the generation pipeline — remains a
        black box in their framework.

  \item \textbf{Tao et al.}~\cite{tao2026ragsurvey} provide the survey most
        thematically adjacent to ours: a focused review of
        \emph{Retrieval-Augmented Code Generation} (RACG) with emphasis on
        repository-level approaches.
        Yet their scope is \emph{one} paradigm (retrieval-augmented generation)
        out of the seven we identify.
        Paradigms based on dense embeddings (PAR-1), structural graphs (PAR-3),
        execution traces (PAR-4), memory augmentation (PAR-5), compression
        (PAR-6), and agentic exploration (PAR-7) fall outside their review
        boundary, leaving the landscape fragmented.

  \item \textbf{Zhang et al.}~\cite{zhang2026surveyse} offer a recent
        comprehensive survey of LLMs for SE in \emph{Science China Information
        Sciences}.
        Their coverage is broad (analogous to Hou et al.) and organised around
        SE tasks and LLM architectures.
        Like Hou et al., the problem of \emph{how} repository context is
        constructed, selected, and compressed receives no dedicated treatment.
\end{itemize}

Table~\ref{tab:survey-comparison} crystallises the positioning of our work
relative to these five surveys across eight analytical dimensions.

\begin{table*}[t]
  \caption{Positioning of this SLR against five closely related surveys.
           \ding{51}~= covered as primary focus; \textbf{(\ding{51})}~=
           partially covered; \ding{55}~= out of scope.}
  \label{tab:survey-comparison}
  \small
  \begin{tabular}{lcccccc}
    \toprule
    \textbf{Dimension} &
    \textbf{Hou 2024}~\cite{hou2024llmse_slr} &
    \textbf{Husein 2025}~\cite{husein2025llmcompletion_slr} &
    \textbf{Jiang 2026}~\cite{jiang2026surveycodegen} &
    \textbf{Tao 2026}~\cite{tao2026ragsurvey} &
    \textbf{Zhang 2026}~\cite{zhang2026surveyse} &
    \textbf{This SLR} \\
    \midrule
    Repository-level scope          & \textbf{(\ding{51})} & \ding{55} & \textbf{(\ding{51})} & \ding{51} & \textbf{(\ding{51})} & \ding{51} \\
    Context representation taxonomy & \ding{55} & \ding{55} & \ding{55} & \textbf{(\ding{51})} & \ding{55} & \ding{51} \\
    7-paradigm RRP taxonomy         & \ding{55} & \ding{55} & \ding{55} & \ding{55} & \ding{55} & \ding{51} \\
    RAG paradigm coverage           & \textbf{(\ding{51})} & \textbf{(\ding{51})} & \textbf{(\ding{51})} & \ding{51} & \textbf{(\ding{51})} & \ding{51} \\
    Agentic paradigm coverage       & \textbf{(\ding{51})} & \ding{55} & \textbf{(\ding{51})} & \textbf{(\ding{51})} & \textbf{(\ding{51})} & \ding{51} \\
    Cross-SE-task mapping (7 tasks) & \ding{51} & \ding{55} & \textbf{(\ding{51})} & \ding{55} & \ding{51} & \ding{51} \\
    10-dimension paradigm analysis  & \ding{55} & \ding{55} & \ding{55} & \ding{55} & \ding{55} & \ding{51} \\
    5-pillar SLR methodology        & \ding{55} & \ding{55} & \ding{55} & \ding{55} & \ding{55} & \ding{51} \\
    \bottomrule
  \end{tabular}
\end{table*}

\subsection{Contributions}
\label{subsec:contributions}

This paper makes the following original contributions:

\begin{enumerate}
  \item \textbf{A unified taxonomy of seven repository representation
        paradigms.}
        We identify, define, and validate seven distinct strategies by which
        LLM-based systems construct repository context:
        (PAR-1) dense code embeddings~\cite{feng2020codebert,guo2022unixcoder},
        (PAR-2) retrieval-augmented context~\cite{lu2022reacc,zhang2023repocoder},
        (PAR-3) graph-based structural representations
        \cite{guo2021graphcodebert,liu2024graphcoder},
        (PAR-4) execution and dynamic representations~\cite{parasaram2023rete,
        wang2023rapgen},
        (PAR-5) memory-augmented context~\cite{wang2026codemem,pan2026prometheus},
        (PAR-6) compression and summarization
        \cite{zan2025codes,shi2025longcodezip},
        and (PAR-7) agentic and iterative exploration
        \cite{yang2024sweagent,bairi2024codeplan,hong2024metagpt}.
        The taxonomy is grounded in 159 primary studies and validated against
        a comprehensive 10-dimension analytical framework covering context
        representation type, granularity, construction pipeline, computational
        complexity, token budget, context freshness, semantic precision,
        scalability, cross-file dependency support, and training dependency.

  \item \textbf{A cross-task mapping of paradigms against 7 SE objectives.}
        We systematically map each of the seven paradigms against seven SE
        tasks: code completion, code generation from natural language,
        automated program repair, fault/bug localization, code search,
        code translation, and test generation.
        This dual-axis analysis (paradigm $\times$ SE task) reveals which
        paradigms dominate which tasks, which paradigms are systematically
        absent from certain tasks, and what open research opportunities exist.

  \item \textbf{A rigorous 5-pillar SLR methodology.}
        The review protocol integrates five complementary standards:
        Kitchenham's SLR guidelines~\cite{kitchenham2007guidelines},
        Wohlin's snowballing procedure~\cite{wohlin2014snowballing},
        PRISMA~2020 reporting standards~\cite{page2021prisma},
        Cruzes and Dyb{\aa}'s thematic synthesis
        methodology~\cite{cruzes2011recommended}, and the
        ACM SIGSOFT Empirical Standards~\cite{acmsigsoft2020standards}.
        This combination ensures \emph{reproducibility}, \emph{transparency},
        and \emph{rigour} at a level not matched by any of the five surveys
        above.
        A full PRISMA 2020 flow diagram and a completed PRISMA checklist
        (Online Supplement) document every stage of the search, screening,
        eligibility assessment, and inclusion process.

  \item \textbf{Empirical evidence of benchmark validity threats.}
        We synthesise evidence that the dominant static-repository assumption
        embedded in PAR-1, PAR-3, and PAR-6 introduces a systematic
        evaluation bias of 10--61\% in benchmark performance estimates
        \cite{zheng2025humanevo}, a threat that all five prior surveys omit.

  \item \textbf{A structured research agenda.}
        We distil 11 open challenges — spanning the temporal freshness problem,
        token-budget optimisation, multi-modal repository representations,
        and evaluation methodology — into a prioritised research agenda for
        the field.
\end{enumerate}

\subsection{Why This Paper, Why Now}
\label{subsec:whynow}

The timing of this review is deliberate.
Between 2024 and early 2026, the repository-level paradigm underwent three
simultaneous transitions:
(i) \emph{agentic systems} (PAR-7) surpassed retrieval-augmented systems
(PAR-2) as the dominant approach on SWE-bench~\cite{jimenez2024swebench,
zhang2025swebenchlive}, with resolution rates climbing from 12\% to over 50\%;
(ii) the \emph{memory} paradigm (PAR-5) emerged as a distinct research thread
\cite{wang2026codemem,chen2025sweexp}, separating persistent cross-session
knowledge from transient in-context retrieval;
and (iii) benchmark construction methodologies matured to expose previously
hidden biases~\cite{zheng2025humanevo,liu2024repobench}, fundamentally
questioning results obtained with earlier evaluation setups.
A survey that predates 2025 — including three of the five surveys reviewed
above — necessarily misses these transitions.
A survey scoped to a single paradigm (Tao et al.) or to a single SE task
(Husein et al.) cannot expose the cross-paradigm trade-offs that practitioners
need to make informed design decisions.
\emph{This SLR is designed to fill precisely that gap.}

\subsection{Paper Organisation}
\label{subsec:organisation}

The remainder of this paper is organised as follows.
Section~\ref{sec:methodology} describes our five-pillar SLR protocol, including
the PRISMA 2020 flow diagram, inclusion/exclusion criteria, quality assessment
procedure, and thematic synthesis method.
Section~\ref{sec:taxonomy} presents the seven-paradigm taxonomy, defines each
paradigm formally, and analyses each paradigm across 10 structural dimensions.
Section~\ref{sec:semapping} maps paradigms against seven SE objectives and
reports cross-paradigm coverage statistics.
Section~\ref{sec:findings} presents the synthesis findings: performance trends,
benchmark threats, open challenges, and emerging hybrid approaches.
Section~\ref{sec:discussion} discusses implications for researchers and
practitioners, limitations of the review, and validity threats.
Section~\ref{sec:relatedwork} positions this work in detail against existing
surveys.
Section~\ref{sec:conclusion} concludes with a structured research agenda.
